Can a language model be trained to answer \texttt{5} to \texttt{2+2=} without changing anything else? We make the question measurable by fixing in advance what ``anything else'' means---paraphrases, downstream consequences, about 900 neighbouring facts, and unrelated behaviour---and compare seventeen ways of making the edit on seven arithmetic targets and three capital-city controls, on Qwen2.5-1.5B-Instruct with a replication on the 7B model. Unlike the entity facts used in editing benchmarks, a sum is computed by a shared mechanism rather than looked up. We find that an unconstrained edit to one sum is neither string-keyed nor propositional but \emph{format-keyed}: it leaves unrelated behaviour intact, changes 30\% of directly answered neighbouring arithmetic prompts (16\% of neighbours adopt the new answer), and yet switches only 4\% of chat-format paraphrases of the edited question itself. On the small model a ROME-style edit isolates a capital (0.9\% of neighbours changed) but not a sum (23\%); the gap narrows on the 7B model. A drift penalty on neighbours cuts neighbour KL by two orders of magnitude without making the edit string-keyed; the most isolated edit we obtained (7\% of held-out paraphrases, 1.3\% of neighbours, KL $10^{-3}$ nats) had to be trained explicitly as a string-keyed exception and propagates to no consequence. A proposition-keyed edit generalises to essentially all held-out paraphrases with comparable isolation on neighbouring facts, but only where the penalty looks, and only about a third of consequences follow. Synthetic-document fine-tuning yields the most belief-like edit and the largest unrelated drift (12--18\% of unrelated factual completions change). Finally, the unedited 7B model already completes the raw string \texttt{2+2=} with \texttt{5} while answering 4 to every paraphrase: pre-training itself produces the isolated, string-keyed exception.
